\documentclass[../main.tex]{subfiles}
\begin{document}
%%%%%%%%%%%%%%%%%%%%%%%%%%%%%%%%%%%%%%%%%%%%%%%%%%%%%%%%%%%%%%%%%%%%%%%%%%%%%%%%%%%%%%%%%%%%%%%%%%%
\chapter{Basic Setup and Configuration}

Everyone has a different style of working, and for that, OmniGraffle offers many ways that you
can customize and configure the way you use the app. This chapter covers how to:

\begin{itemize}
  \item Choose or create a default template that OmniGraffle uses whenever you create a new
    document
  \item Set options in the Preference panes (\textbf{Command-,}) to make your workflow smoother
  \item Use the new \textbf{Keyboard Shortcuts} window to change OmniGraffle's default
    shortcuts, or share shortcut sets with other users
  \item Change the default set of buttons you see in toolbar
  \item Rearrange the tools found in the \textbf{Tool Palette}
  \item Rearrange the inspectors in the Inspector sidebar
  \item ``Tear-off'' OmniGraffle's inspectors to create custom \textbf{Workspaces} to give you
    quick access to \textit{just} the inspectors you need
\end{itemize}

Part FAQ, part tutorial, this chapter provides you with the assistance you need to customize
your OmniGraffle user experience.

\section{Choosing a Default Template}

One of the most frequent questions our Support Humans receive is \textit{``How do I change the
default template in OmniGraffle?''}, so let's get this one out of the way first.

The first way to change OmniGraffle's default template is in the \textbf{General Preferences}
pane (\textbf{OmniGraffle} $\blacktriangleright$ \textbf{Preferences} $\blacktriangleright$
\textbf{General}):

\ogfigure{0.800}{res/og7_config01.png}

\begin{enumerate}
  \item Next to \textbf{New Documents}, click \textbf{Choose}
  \item A sheet slides down, providing you with access to OmniGraffle's available templates.
  \item Click to select the desired template, and then click \textbf{Set as Default}.
\end{enumerate}

\ogfigure{0.800}{res/og7_config02.png}

The next time you create a New Document, OmniGraffle will use your choice of template as the
default.

The next way to set a default template in OmniGraffle is with the \textbf{Resource Browser}. To
open the Resource Browser, choose \textbf{File} $\blacktriangleright$ \textbf{Resource
Browser} (\textbf{Shift-Command-N}):

\ogfigure{0.800}{res/og7_config03.png}

\begin{enumerate}
  \item In the Resource Browser's sidebar, select one of the folders under \textbf{Templates},
    and then browse around until you find a template you'd like to set as the default.
  \item Click to select the desired template.
  \item In the Resource Browser's toolbar, click
    \ogicon{12.5pt}{res/og7_button_rb_actionmenu.png}, and then choose \textbf{Set as Default}
    from the Action Menu.
\end{enumerate}

\ogfigure{0.800}{res/og7_config04.png}

\begin{ognote}
  To learn how to create a new template based on your own needs, see Creating a New Template,
  later in this guide.
\end{ognote}

If you ever need to switch OmniGraffle back to using the default \textbf{Blank} template, go to
\textbf{OmniGraffle} $\blacktriangleright$ \textbf{Preferences} $\blacktriangleright$
\textbf{General}, and then click the \textbf{Reset} button in the General preference pane.

\section{Adding Resources to OmniGraffle}

If you ever need to add a new stencil or template file to OmniGraffle, the easiest way to do
this is with the Resource Browser:

\begin{enumerate}
  \item Choose \textbf{File} $\blacktriangleright$ \textbf{Resource Browser}
    (\textbf{Shift-Command-N}).
  \item In the Resource Browser's toolbar, click
    \ogicon{12.5pt}{res/og7_button_rb_actionmenu.png}, and then choose \textbf{Add Resources to
    Library} from the Action Menu:
\end{enumerate}

\ogfigure{0.800}{res/og7_addresources.png}

OmniGraffle opens a window that you can use to find, select, and add OmniGraffle Stencil and
Template files to make them available in OmniGraffle 7. Use the window's \textbf{Search} field
to search for stencil and template file extensions (\textit{.gstencil} and
\textit{.gtemplate}, respectively).

\ogfigure{0.800}{res/og7_addresource_search.png}

Command-click the files you would like to import, and then click \textbf{Open}.

\begin{ognote}
  See Importing Microsoft Visio Files for information on how to open and export Microsoft Visio
  files with OmniGraffle Pro.
\end{ognote}

\section{Setting OmniGraffle's Preferences}

The most common way to configure any Mac app is by changing some of the default
\textit{preferences}. To open OmniGraffle's preferences window, choose \textbf{OmniGraffle}
$\blacktriangleright$ \textbf{Preferences}, or use the standard \textbf{Command-,} keyboard
shortcut.

Two features that all of the OmniGraffle preference panes offer are the ability to
\textbf{Reset} (\ogicon{12.5pt}{res/og7_button_reset.png}) any changes you've made to their
default state, and quick access to \textbf{Help} (\ogicon{12.5pt}{res/og7_button_help.png})
documentation for that pane.

The following sections closely examine all of the preferences you'll find in OmniGraffle.

\subsection{General Preferences}

The General preferences give you control over basic OmniGraffle settings, such as what
OmniGraffle does at startup, whether new documents open with a blank canvas or you're taken to
the template chooser, and things like text editing and selection behaviors. You can also turn
on \textbf{Multi-Touch} for trackpad support, and opt for OmniGraffle to only show you the
inspectors you need when you're working on the canvas.

\ogfigure{0.800}{res/og7_prefs_general_standard.png}

Options include:

\begin{description}[style=nextline]
  \item[{Startup}] Select \textbf{Create new document if nothing else is open} to start with a
    new document when you aren't making OmniGraffle remember open documents (with \textbf{System
      Preferences} $\blacktriangleright$ \textbf{General} $\blacktriangleright$ \textbf{Close
    windows when quitting an app}), or if there were no open documents the last time you quit.
  \item[{New Documents}] Select \textbf{Open Resource Browser} to be presented with a window for
    choosing from a variety of file sources each time you create a new document. Select
    \textbf{Use Template}, and choose your default template, if you'd rather bypass the browser.
    Either way, you can always open the Resource Browser from the \textbf{File} menu or with the
    \textbf{Shift-Command-N} keyboard shortcut.
  \item[{Help Tags}] When you move or resize an object, a little help tag appears to let you
    know the object's exact new location or size.
\end{description}

\begin{ogpro}
  \begin{description}[style=nextline]
    \item[{Help Tags (\textbf{PRO})}] If you have OmniGraffle Pro, a second option is available
      for displaying any notes and metadata you've assigned to an object in the Object Data
      inspector when you mouse over them.
  \end{description}
\end{ogpro}

\begin{ogpro}
  \begin{description}[style=nextline]
    \item[{Object ID Numbers (\textbf{PRO})}] If you're accessing OmniGraffle objects with
      AppleScript, or you're otherwise endeared to unique IDs, you can select \textbf{Show
      object identification numbers to add a column to the list sidebar}.
  \end{description}
\end{ogpro}

\begin{description}[style=nextline]
  \item[{Full Screen}] Choose whether to display the toolbar when in full screen mode.
  \item[{Inspectors}] Choose whether to hide inspectors that don't apply to the currently
    selected object. You can also opt to turn on the ability to define stroke widths based on
    the template's \textbf{Ruler Units} instead of the default, \textbf{points}.
  \item[{Stencil Browser}] By default, OmniGraffle's Stencil Browser is available in the right
    sidebar. However, you can use the \textbf{Stencil Browser} pop-up menu to change the default
    location to one of the following:

    \begin{itemize}
      \item Left Sidebar
      \item Right Sidebar
      \item Popover
      \item Floating Window
    \end{itemize}
  \item[{Sidebar}] Choose whether you'd like to edit the names of newly-created canvases and
    layers in the Sidebar. When new canvases and layers are added to a document, their default
    names (e.g., Canvas 1, Canvas 2, or Layer 1, Layer 2, and so on) are highlighted in the
    Sidebar so you can enter a relevant name.
\end{description}

\subsection{Drawing Tools}

The Drawing Tools preferences provide you with a way to organize the tools you see in
OmniGraffle's toolbar. You can also define how you interact with the tools, and set the
behaviors for line creation and editing.

\ogfigure{0.800}{res/og7_prefs_drawingtools_standard.png}

\begin{description}[style=nextline]
  \item[{Keep tool active}] By default, it takes one click to temporarily activate a tool, and
    two clicks to make a tool ``stick'' in an active state. You can choose to make the tools
    stick after only one click, or make it so that every tool selection is temporary.
  \item[{Line Creation}] When you are creating a line, clicking an object will always connect
    the line to the object and end the line. If you click in an empty area of the canvas,
    though, you might want to end the line there or you might want to add a point there and
    continue creating the line. Use this preference setting to indicate which you prefer.
  \item[{Line Editing}] When you double-click a line, you might prefer to add a new point to the
    line, or you might prefer to add a label. Choose one; \textbf{Option-double-click} performs
    the opposite action.
\end{description}

Clicking the \textbf{Customize Keyboard Shortcuts} button opens the \textbf{Keyboard Shortcuts}
window, which you can use to set, change, or remove any keyboard shortcut in OmniGraffle. Each
menu item is represented in the Keyboard Shortcuts window, allowing you to set shortcuts for
submenu and sub-submenu items as well.

\subsubsection{Customizing the Tool Palette}

You can change which tools are available on the left side of the Tool Palette by moving the
dividing bar or dragging the tools into a different position.

\ogfigure{0.673}{res/og7_setup_toolpalette_rearrange.png}

For example, if you mostly use the Selection, Shape, Pen, and Point Editor tools, you can could
set up your toolbar as follows:

\ogfigure{0.673}{res/og7_setup_toolpalette.png}

Which results in the Tool Palette appearing as follows when collapsed:

\ogfigure{0.673}{res/og7_setup_toolpalette_toolbar.png}

\subsection{Canvas}

The Canvas preferences pertain to how you work with things on the Canvas.

\ogfigure{0.800}{res/og7_prefs_canvas_standard.png}

\begin{description}[style=nextline]
  \item[{Text Editing}] Decide how you would like the \textbf{Return} and \textbf{Tab} keys to
    behave when you are editing the label of an object or line, or when working with text blocks
    on the canvas:

    \begin{itemize}
      \item \textbf{Return} --- When pressing \textbf{Return}, you can opt to have that start a
        new line of text or to end editing (similar to pressing the \textbf{Escape} key).
      \item \textbf{Tab} --- When pressing \textbf{Tab}, you can opt to have that insert a Tab
        character, or to switch from the currently selected object to the next object so you can
        edit its text.
    \end{itemize}

    Whichever way you choose, you can hold \textbf{Option} while pressing \textbf{Return} or
    \textbf{Tab} to do the reverse instead.
  \item[{Selection}] These two options further define the behavior of selected objects:

    \begin{itemize}
      \item \textbf{Click through objects with no fill} --- objects that don't have a fill
        applied, such as \textbf{Artboard objects}, can be clicked through regardless of their
        placement in the layer stack.
      \item \textbf{Change active layer upon selection} --- when you click to select an object,
        that object will be selected, regardless of which layer the object is on. If you
        frequently work on projects with multiple layers, you should consider enabling this
        option.
    \end{itemize}
  \item[{Alignment}] Decide which object to use as the default anchor when aligning objects on
    the canvas (using the Alignment inspector or Arrange menu, for example):

    \begin{itemize}
      \item \textbf{Align to selection} --- the object in the selection that is nearest to the
        desired alignment position is used as the anchor (other objects are aligned to it).
      \item \textbf{Align to first-selected object} --- when selecting multiple objects
        consecutively, other objects are aligned to the first one selected.
    \end{itemize}
\end{description}

\subsection{Multi-Touch}

If you are using a laptop or touch-sensitive trackpad, you can enable Multi-Touch trackpad
gestures. For example, you can use the two-finger touch gesture to zoom in or out on the
Canvas, or to rotate a selected object.

\ogfigure{0.800}{res/og7_prefs_multitouch_standard.png}

OmniGraffle can also provide \textbf{Haptic Feedback} through the trackpad (on supported
devices) for things such as guides, object rotation, and whenever you use a stepper control in
one of the inspectors.

\subsection{Appearance}

The Appearance preference panel allows you to change the default colors used for Ruler Guides,
Smart Guides, and Bézier control handles. You can also change the size of Bézier control
points (this is a visual change only, for ease of interaction with the points on the canvas).

\ogfigure{0.800}{res/og7_prefs_appearance_standard.png}

You can also change the colors of Ruler Guides in the Sidebar (\textbf{Option-Command-3}). For
more information on Ruler Guides, see The Guides Tab in Managing Canvases, Layers, and Objects
in the Sidebar.

\subsection{Presentation (\texorpdfstring{\textbf{PRO}}{PRO})}

\begin{ogpro}
  For those times when you need to make a presentation, say to show off the new user interface for
  an app you're developing or when you're showing a client how you plan to redesign their home,
  OmniGraffle Pro offers Presentation mode. Available by choosing \textbf{View}
  $\blacktriangleright$ \textbf{Start Presentation}, OmniGraffle takes over your entire screen
  so you can easily show off your work to everyone in the room.

  Use the Presentation preferences to configure your settings for making presentations.

  \ogfigure{0.800}{res/og7_prefs_presentation_pro.png}

  \begin{description}[style=nextline]
    \item[{Highlight Behavior}] Choose whether to have an object become highlighted whenever you
      put the pointer on it, whenever you click it, or never. You can also choose to display a
      badge for objects that have actions set up in the Action inspector.
    \item[{Highlight Appearance}] Click the \textbf{Color} well
      \ogicon{12.5pt}{res/og7_prefs_highlightwidget.png} to choose a color for the highlight.
      Use the slider to adjust the thickness of the highlight.
  \end{description}

  To learn more about using OmniGraffle's Presentation Mode, see Using OmniGraffle in
  Presentation Mode (PRO) later in this guide.
\end{ogpro}

\subsection{Update Preferences}

If you've purchased your copy of OmniGraffle 7 from the Omni Group's website, OmniGraffle can
automatically use your Internet connection to check for updated releases.

\begin{ognote}
  The Update Preferences panel doesn't exist if you have purchased OmniGraffle 7 from the Mac App
  Store. When a new version of OmniGraffle is available, the Mac App Store will alert you to
  download and install that on your Mac.
\end{ognote}

If \textbf{Check for updates} is turned on, you will be automatically notified when there is a
new version of OmniGraffle ready for you to download and enjoy. Click \textbf{Check for
Updates Now} to look for the newest version manually; you can also choose \textbf{OmniGraffle}
$\blacktriangleright$ \textbf{Check for Updates} from the menu bar. This is particularly
helpful if you have \textbf{Check for updates} set to check Monthly and you've recently heard
that a new update is available.

\ogfigure{0.800}{res/og7_prefs_update_standard.png}

Additionally, you can opt in to \textbf{Include system information} when OmniGraffle checks for
updates. This sends basic information about your system to us, which we \textit{only} use to
better support various systems with future releases of OmniGraffle. And, as the Update
preferences pane says: \textit{``This information is kept entirely anonymous,''} which means
that the information we receive about your system stays with us and is only used by us as we
work on future updates of OmniGraffle.

\section{Customizing OmniGraffle's Keyboard Shortcuts}

Keyboard shortcuts make it easy for you to quickly access menu options without having to use a
mouse. OmniGraffle has an extensive set of default keyboard shortcuts, which you can change by
choosing \textbf{OmniGraffle} $\blacktriangleright$ \textbf{Keyboard Shortcuts} from the menu
bar.

\ogfigure{0.800}{res/og7_shortcuts01.png}

The Keyboard Shortcuts window is divided into two sections: \textbf{Tools} and \textbf{Menus}:

\begin{itemize}
  \item \textbf{Tools} --- lists the tools found in the \textbf{Tool Palette} at center of the
    toolbar; things you use in OmniGraffle to create lines, shapes, and such.
  \item \textbf{Menus} --- lists the menus found at the top of the screen, along with each menu
    item found in that menu or submenu.
\end{itemize}

You can hide and show each section by clicking to the right of those sections (the \textbf{Hide}
button text appears on mouseover). You can also expand and collapse the menu subsections by
clicking the disclosure triangles (\ogicon{8.3pt}{res/og7_disclosure_closed.png}
\ogiconw{7.5pt}{res/og7_disclosure_open.png}) along the left.

Keyboard shortcuts for Tools can either be a single letter or number without the use of special
key combinations, such as the \textbf{Control}, \textbf{Option}, \textbf{Shift}, or
\textbf{Command} keys. This makes it really easy for you to access and use the Tool you need
with just a single key press.

\begin{ogtip}
  With keyboard access to OmniGraffle's tools, a single press of the keyboard shortcut (for
  example, pressing the \textbf{s} key for the \textbf{Shape} tool) activates the tool.

  To make the tool active, you can either press and hold the tool's shortcut key, or press the
  key a second time to make it sticky.
\end{ogtip}

Use the \textbf{Search} field to search for a Tool or Menu option. This makes it easy for you to
see how the assigned keyboard shortcuts relate to each other, and to reassign or set keys, if
desired.

\ogfigure{0.800}{res/og7_shortcuts02.png}

\subsection{Editing OmniGraffle's Keyboard Shortcuts}

Changing a keyboard shortcut is easy; simply click the field to the right of a Tool or Menu item
and enter a new shortcut. The changes are automatically available for use, without requiring
you to quit and restart OmniGraffle.

If you enter a shortcut that is used by another Tool or Menu item, the Keyboard Shortcuts window
displays a warning symbol within both shortcut fields, and the normally-grayed out warning
button at the top of the window becomes available. This makes it easy for you to identify, and
correct, the conflict.

\ogfigure{0.800}{res/og7_shortcuts03.png}

Pressing the Warning button at the top of the Keyboard Shortcuts window focuses on the
conflicting items, and a \textit{revert} arrow \ogicon{12.5pt}{res/og7_icon_revertarrow.png}
is placed to the right of the item that's causing the conflict. Clicking
\ogicon{12.5pt}{res/og7_icon_revertarrow.png} changes that particular item back to its default
keyboard shortcut, or no shortcut if none was previously assigned.

To remove a keyboard shortcut for an item, click the field to the right of a Tool or Menu item
and press the \textbf{Delete} key. This removes the keyboard shortcut for that item, which can
be useful if the \textit{revert} arrow sets the key binding to a conflicting value.

The moment you make a change to the Default keyboard set, the name of the current keyboard set
used in the drop-down menu becomes editable. By default, OmniGraffle places your user account
name in the field, but you can change this to whatever you'd like; for example, \textbf{Design
Set} or \textbf{Diagramming}.

\ogfigure{0.800}{res/og7_shortcuts04.png}

If you later decide that you want to change the name of a keyboard set, do the following:

\begin{enumerate}
  \item Click \ogicon{12.5pt}{res/og7_button_actionmenu.png} next to the drop-down menu and
    choose the \textbf{Rename} option from the menu:

    \ogfigure{0.800}{res/og7_shortcuts05.png}
  \item The name of the keyboard set in the drop-down menu becomes editable. Enter a new name,
    and press \textbf{Return} to apply the new name.

    \ogfigure{0.800}{res/og7_shortcuts06.png}
\end{enumerate}

\subsection{Using Another Keyboard Set}

Because we know designers have many tools in their toolbox, OmniGraffle 7 also ships with
keyboard shortcut sets based on Adobe Illustrator\textsuperscript{\textregistered{}} and
Sketch. If you use one of those apps alongside OmniGraffle, you can use the drop-down menu at
the top of the Keyboard Shortcuts window to choose the set that most closely matches your
workflow in other applications.

\ogfigure{0.673}{res/og7_shortcuts07.png}

To revert the shortcuts back to OmniGraffle's default set, choose \textbf{OmniGraffle}
$\blacktriangleright$ \textbf{Keyboard Shortcuts} and then select \textbf{Default} in the
drop-down menu.

\subsection{Sharing Keyboard Sets (\texorpdfstring{\textbf{PRO}}{PRO})}

\begin{ogpro}
  After changing some keys and setting the name for your new keyboard set, you can export your
  keyboard set by clicking \ogicon{12.5pt}{res/og7_button_actionmenu.png} and choosing
  \textbf{Export Shortcuts}. When you export a keyboard set, OmniGraffle uses the name of the
  set by default, paired with a \textit{.graffle-shortcuts} file extension.

  \ogfigure{0.800}{res/og7_shortcuts08.png}

  \begin{ogwarning}
    You can change the first part---the filename---but \textit{do not} change the file extension.
    OmniGraffle expects the \textit{.graffle-shortcuts} file extension on keyboard sets you
    import.
  \end{ogwarning}

  Similarly, you can also import a keyboard set, using one of the following methods:

  \begin{itemize}
    \item Click \ogicon{12.5pt}{res/og7_button_actionmenu.png} and choose \textbf{Import
      Shortcuts}. In the sheet that slides down, locate and select the
      \textit{.graffle-shortcuts} file you'd like to import, and then click \textbf{Open} to add
      that keyboard set to the Keyboard Shortcuts drop-down menu.
    \item Use the Finder to locate and then double-click the \textit{.graffle-shortcuts} file.
      The Keyboard Shortcuts window will open with that shortcuts file already selected in the
      pop-up menu.
  \end{itemize}
\end{ogpro}

\section{Customizing the Toolbar}

Along the top of OmniGraffle's window is the \textbf{Toolbar}, packed with buttons that can be
customized to suit the needs of your workflow and most common interactions.

\ogfigure{0.800}{res/og7_toolbar01.png}

Like most Toolbars, OmniGraffle's comes with a default set of buttons, based on what we believe
are common things you'd want to do. However, you can customize the toolbar by adding
additional buttons or rearranging their order, to suit your needs.

To customize the toolbar, choose \textbf{View} $\blacktriangleright$ \textbf{Customize Toolbar}
or \textbf{Control-click} on the toolbar and select \textbf{Customize Toolbar}. When you do, a
sheet slides down from under the toolbar, revealing a slew of function-related buttons that
you can drag up to the toolbar. For example, you could add the \textbf{New Artboard Layer} or
\textbf{Start Presentation} buttons to the toolbar if you've purchased OmniGraffle 7 Pro.

\ogfigure{0.800}{res/og7_toolbar02.png}

If you don't like the changes you've made, or if you ever want to reset the toolbar to it's
default state, all you need to do is drag the bottom row up to the toolbar.

Click \textbf{Done} when you've finished making changes to the toolbar.
%%%%%%%%%%%%%%%%%%%%%%%%%%%%%%%%%%%%%%%%%%%%%%%%%%%%%%%%%%%%%%%%%%%%%%%%%%%%%%%%%%%%%%%%%%%%%%%%%%%
\end{document}
